% Created (UTC): 2026-06-25T21:43:04Z
% Repository HEAD: 467d9018ec87aea7ac8a5550a166386d95ae6cdb
\documentclass[11pt]{article}
\usepackage[margin=1in]{geometry}
\usepackage{amsmath,amssymb,amsthm,mathtools}
\usepackage{booktabs}
\usepackage{microtype}
\usepackage{enumitem}
\emergencystretch=3em \hbadness=4000
\usepackage[colorlinks=true,linkcolor=blue,citecolor=blue,urlcolor=blue]{hyperref}
\DeclareMathOperator{\Li}{Li}
\DeclareMathOperator{\Real}{Re}
\newcommand{\Cl}{\operatorname{Cl}_2}
\newcommand{\Jcal}{\mathcal{J}}
\newcommand{\eps}{\varepsilon}
\newcommand{\rad}{\operatorname{rad}}
\theoremstyle{plain}\newtheorem{theorem}{Theorem}\newtheorem{proposition}[theorem]{Proposition}
\newtheorem{corollary}[theorem]{Corollary}
\theoremstyle{remark}\newtheorem{remark}[theorem]{Remark}
\title{Arithmetic of the Herglotz function:\\
verified tables, new rational-value theorems, and bridges to the\\ golden ladders and the Clausen lattice}
\author{Vladimir Reshetnikov}
\date{2026-06-25}
\begin{document}
\maketitle
\begin{center}\small
Consolidated research report. Created (UTC) 2026-06-25T21:43:04Z; repository \texttt{HEAD}
\texttt{467d9018ec87aea7ac8a5550a166386d95ae6cdb}.\\
Source paper: D.\ Radchenko and D.\ Zagier, \emph{Arithmetic properties of the Herglotz function}
(\texttt{src/PolyLog/docs/pdfs/HerglotzFunction/}).\\
Stores: \texttt{herglotz-rational-values.wl}, \texttt{herglotz-quadratic-units.wl}.
Reproducers: \texttt{src/PolyLog/tools/scratch/herglotz/}.
\end{center}

\begin{abstract}
This report consolidates an experimental-mathematics study of the Herglotz function $F(x)$, building on
Radchenko and Zagier's paper. We do three things. First, we re-implement $F$, its companion $\Jcal$, and the
exact Theorem-3 dilogarithm sum independently (PARI/GP and \texttt{mpmath}) and re-verify \emph{every} entry
of the authors' Tables~1 and~2 of quadratic-unit values to better than $10^{-80}$; two rows of Table~1 fail,
and we identify them as genuine typographical errors in the printed paper, each fixed by a unique single
edit. Second, we turn the paper's two worked rational-argument examples into a family of general theorems:
a closed form for $F(2/q)$ valid for all $q$ (with proof), the pure-logarithm $F(1/q)$, the explicit
$F(n/(n+1))$ family, a reciprocity law for $F(p/2)$, a derivative law for $F'(2/q)$, a clean collapse
criterion ($F(p/q)$ reduces to rational-argument dilogarithms iff $\rad(q)\mid 30$), and the determination
that the paper's pure-logarithm value $J(2/5)$ is special to $q\in\{3,5\}$. Third, we connect the new
constants to two adjacent corpora: the golden-argument dilogarithms of $F(2/5)$ to the golden polylogarithm
ladders over $\mathbb Q(\sqrt5)$, and the Clausen sum of the derivative law to the conductor-$q$ Clausen
lattice at roots of unity. Every displayed identity is verified numerically; a residual table is given in
\S\ref{sec:verif}.
\end{abstract}

\tableofcontents

\section{The function and the exact dilogarithm sum}
The Herglotz function may be taken in the integral form (the source paper's eq.~(1))
\begin{equation}\label{eq:Jint}
  J(x)=\int_0^1\frac{\log(1+t^{x})}{1+t}\,dt,\qquad J(x)=F(2x)-2F(x)+F(x/2)+\frac{\pi^2}{12x},
\end{equation}
which relates the companion $J$ to the Herglotz function $F$. For the rational-argument results the engine is
Radchenko--Zagier's Theorem~3, an \emph{exact, finite} evaluation obtained from the second integral
representation of $F$ via the distribution relations:
\begin{equation}\label{eq:dilogsum}
  F\!\Bigl(\tfrac PQ\Bigr)-F(1)=\sum_{\alpha^P=1}\Bigl(\sum_{\substack{\beta^Q=1\\\beta\neq1}}f(\alpha,\beta)
  -\sum_{\substack{\beta^P=1\\\beta\neq1}}f(\alpha,\beta)\Bigr),\quad
  f(\alpha,\beta)=\Li_2\!\Bigl(\tfrac{\beta}{\beta-1}\Bigr)-\Li_2\!\Bigl(\tfrac{\alpha-\beta}{1-\beta}\Bigr).
\end{equation}
Since~\eqref{eq:dilogsum} is a finite sum of dilogarithms it evaluates $F(P/Q)-F(1)$ to arbitrary precision
with no quadrature error; throughout we use it as ground truth and cross-check against the defining integral
(which is independent but only $\sim\!30$-digit accurate through \texttt{mpmath} quadrature).

\section{Quadratic-unit values: Tables 1 \& 2, and two errata}\label{sec:tables}
For real quadratic units the relevant object is the companion
$\Jcal(x)=J(x)-\tfrac{\log^2 2}{2}+\tfrac{\pi^2}{24}(x-\tfrac1x)$, and the arithmetic atoms are
$S_{p,q}=\log\eps_p\,\log\eps_q$ where $\eps_d$ is the fundamental unit of $\mathbb Q(\sqrt d)$ for a
fundamental discriminant $d>1$ and $\eps_1:=2$. Tables~1 and~2 of the source paper express, for the $45$
known one-class-per-genus cases, the quantities
$4\bigl(\Jcal(n+\sqrt{n^2-1})-n\tfrac{\pi^2}{24}\bigr)$ and
$2\bigl(\Jcal(n+\sqrt{n^2+1})-n\tfrac{\pi^2}{24}\bigr)$ as integer combinations of the $S_{p,q}$.

We re-implemented $\Jcal$ and $S_{p,q}$ in PARI/GP (\texttt{intnum} integrals, \texttt{bnfinit} regulators)
and evaluated both sides of all $45$ rows at $200$-digit precision. All $16$ rows of Table~2 and $27$ of the
$29$ rows of Table~1 agree to better than $10^{-80}$. The two failures are $O(1)$:
\[
  \text{$n=43$:}\ \ \text{LHS}-\text{RHS}=-13.0160\ldots,\qquad
  \text{$n=61$:}\ \ \text{LHS}-\text{RHS}=-2.9011\ldots .
\]
Each resolves to a single exact correction, confirmed against the high-DPI scan to be a typo in the printed
paper rather than an OCR artifact.

\begin{proposition}[Two corrections to Table 1]
\leavevmode
\begin{itemize}[topsep=2pt,itemsep=1pt]
\item $n=43$ ($n^2-1=2^3\cdot3\cdot7\cdot11$): the printed right-hand side overshoots by exactly $S_{33,56}$
$\bigl((\text{LHS}-\text{RHS})/S_{33,56}=-1.000\ldots$ to $110$ digits$\bigr)$; the spurious $+S_{33,56}$
term should be deleted.
\item $n=61$ ($n^2-1=2^3\cdot3\cdot5\cdot31$): the printed $S_{60,1032}$ should read $S_{60,248}$. The
printed subscript $1032=2^3\cdot3\cdot43$ carries the prime $43$, foreign to $n^2-1=3720$, whereas
$60\cdot248=14880=4\cdot3720$ is the correct genus splitting.
\end{itemize}
A brute-force search over all single edits (delete a term, change a coefficient in $[-4,4]$, or replace a
subscript by any admissible fundamental-discriminant pair) returns these as the \emph{unique} fixes. With
both applied, all $45$ rows verify to better than $10^{-80}$.
\end{proposition}

\noindent A by-product of the same study maps where $\Jcal$ ceases to be elementary: $\Jcal(x)$ closes in
$\{\zeta(2),\log2\log x,\log^2 x\}$ precisely for the even-trace family $x=n+\sqrt{n^2\pm1}$
(e.g.\ $\Jcal(1+\sqrt2)=\tfrac14\zeta(2)+\tfrac12\log2\log(1+\sqrt2)$), and \emph{not} for a generic
quadratic unit such as the golden ratio --- the precise numerical content of the paper's remark that the
individual values $F(x)$ have no general closed form.

\section{Rational arguments: a family of theorems}\label{sec:rational}
The paper writes out only two rational-argument cases: $F(n)$ (pure logarithms, its eq.~23) and the lone
fractional example $F(2/5)$ (its eq.~24). We now give the general laws.

\subsection{Two reductions on the line \texorpdfstring{$\Real=\tfrac12$}{Re=1/2}}
Everything rests on one observation: for a unit $\beta=e^{i\theta}$,
\begin{equation}\label{eq:line}
  \frac{\beta}{\beta-1}\ \text{has real part}\ \tfrac12,\qquad
  \arg\frac{\beta}{\beta-1}=\frac{\theta}{2}-\frac{\pi}{2},\qquad
  \Bigl|\frac{\beta}{\beta-1}\Bigr|=\frac{1}{2\sin(\theta/2)} .
\end{equation}
By the reflection $\Li_2(z)+\Li_2(1-z)=\tfrac{\pi^2}{6}-\log z\log(1-z)$ with $1-z=\bar z$ on $\Real z=\tfrac12$,
\begin{equation}\label{eq:reRe}
  2\Real\Li_2\Bigl(\tfrac12+iy\Bigr)=\frac{\pi^2}{6}-\log^2|z|-(\arg z)^2 .
\end{equation}
The second reduction handles the $\alpha=-1$ part of~\eqref{eq:dilogsum}: $\tfrac{-1-\beta}{1-\beta}=-i\cot\tfrac\theta2$
is purely imaginary, and for real $y$
\begin{equation}\label{eq:imax}
  \Real\Li_2(iy)=\tfrac14\Li_2(-y^2).
\end{equation}

\subsection{$F(1/q)$: pure logarithms}
\begin{theorem}\label{thm:F1q}
For all $q\ge2$, with $m=\lfloor(q-1)/2\rfloor$,
\[
  F\!\Bigl(\tfrac1q\Bigr)-F(1)=-\frac{(q-1)(3q-2)}{24q}\,\pi^2
  -\sum_{k=1}^{m}\log^2\!\Bigl(2\sin\tfrac{\pi k}{q}\Bigr)\;-\;[\,q\ \mathrm{even}\,]\,\tfrac12\log^2 2 .
\]
\end{theorem}
\begin{proof}
With $P=1$ only $\alpha=1$ occurs and the inner $\beta^P$-sum is empty, so~\eqref{eq:dilogsum} reads
$F(1/q)-F(1)=\sum_{\beta^q=1,\beta\neq1}(\Li_2(\tfrac{\beta}{\beta-1})-\tfrac{\pi^2}{6})$. Pairing $\beta$ with
$\bar\beta$ and applying~\eqref{eq:reRe} gives the $\pi^2$ coefficient
$-\bigl[\tfrac{q-1}{12}+\sum_{k=1}^m(\tfrac kq-\tfrac12)^2\bigr]=-\tfrac{(q-1)(3q-2)}{24q}$ and the log-sine sum.
For even $q$ the self-conjugate root $\beta=-1$ contributes $\Li_2(\tfrac12)=\tfrac{\pi^2}{12}-\tfrac12\log^2 2$.
\end{proof}

\subsection{$F(2/q)$: the general theorem}
\begin{theorem}\label{thm:F2q}
For every $q\ge3$,
\[
  F\!\Bigl(\tfrac2q\Bigr)-F(1)=-\frac{(q-1)(q-2)}{12q}\pi^2+\log^2 2
  -2\!\!\sum_{k=1}^{\lfloor(q-1)/2\rfloor}\!\!\log^2\!\Bigl(2\sin\tfrac{\pi k}{q}\Bigr)
  -\tfrac12\!\!\sum_{k=1}^{\lfloor(q-1)/2\rfloor}\!\!\Li_2\!\Bigl(-\cot^2\tfrac{\pi k}{q}\Bigr),
\]
where for even $q$ the two sums run to $k=\tfrac q2-1$ and a further $-\log^2 2$ is added.
\end{theorem}
\begin{proof}[Proof (odd $q$)]
Here $\alpha\in\{1,-1\}$ and the inner $\beta^P$-sum is the single term $\beta=-1$. For the first part,
$f(1,\beta)+f(-1,\beta)=2\Li_2(\tfrac{\beta}{\beta-1})-\tfrac{\pi^2}{6}-\Li_2(\tfrac{-1-\beta}{1-\beta})$.
Summing over $\beta\neq1$ and using~\eqref{eq:reRe} for the first pieces and~\eqref{eq:imax}
(with $\tfrac{-1-\beta}{1-\beta}=-i\cot\tfrac{\pi k}q$, pairing $k\leftrightarrow q-k$) for the last gives
$2S_1-(q-1)\tfrac{\pi^2}{6}-\tfrac12\sum_{k}\Li_2(-\cot^2\tfrac{\pi k}q)$, where
$S_1=\sum_{k=1}^{(q-1)/2}[\tfrac{\pi^2}{6}-\log^2(2\sin\tfrac{\pi k}q)-\pi^2(\tfrac kq-\tfrac12)^2]$. The boundary
gives $-\sum_\alpha f(\alpha,-1)=+\log^2 2$ from $\Li_2(\tfrac12)$. Collecting, the $\tfrac{\pi^2}{6}$ terms
cancel and the $\pi^2$ coefficient is $-2\sum_{k=1}^{(q-1)/2}(\tfrac kq-\tfrac12)^2$, which equals
$-\tfrac{(q-1)(q-2)}{12q}$ by an elementary identity (verified symbolically). The even-$q$ case differs only
in that $\beta=-1$ is now a genuine $q$-th root, contributing the extra $-\log^2 2$.
\end{proof}
The coefficient is exactly twice that of eq.~23 for $F(q)$. The new ingredient $\Li_2(-\cot^2\frac{\pi k}q)$
is a dilogarithm at an algebraic argument of degree $\varphi(q)/2$: for $q=5$ the two terms collapse (the real
subfield is quadratic) to the single $\Li_2(4/5)$ of eq.~24, but for $q\ge7$ they do not, which is why no
simpler form than Theorem~3 exists there. A direct PSLQ search shows the sum
$\sum_k\Li_2(-\cot^2\frac{\pi k}q)$ is itself irreducible over $\{\pi^2,\log^2 q,\text{Catalan},\log\text{-sines}\}$.

\subsection{The $F(n/(n+1))$ family and the $F(p/2)$ reciprocity}
\begin{theorem}\label{thm:np1}
For all $n\ge2$, with $S(q)=\sum_{k=1}^{q-1}\log^2(2\sin\tfrac{\pi k}q)$,
\[
  F\!\Bigl(\tfrac{n}{n+1}\Bigr)-F(1)=-\frac{3n^2+3n+2}{24n(n+1)}\pi^2+\Li_2\!\Bigl(\tfrac{n}{n+1}\Bigr)
  +\tfrac12\log^2\!\tfrac{n+1}{n}-\tfrac12 S(n+1)+\tfrac12 S(n),
\]
the single dilogarithm $\Li_2(n/(n+1))$ appearing with coefficient exactly $+1$ (the source paper's
``bilinear-log'' Example~2, made explicit). The $\pi^2$ coefficient equals $-\tfrac18-\tfrac1{12n(n+1)}$.
\end{theorem}
\begin{theorem}[Reciprocity]\label{thm:p2}
For all odd $p\ge3$,
\[
  F\!\Bigl(\tfrac p2\Bigr)+F\!\Bigl(\tfrac2p\Bigr)-2F(1)=-\frac{(p-2)^2}{12p}\pi^2+\tfrac12\log^2\!\tfrac p2 .
\]
The $\Li_2(-\cot^2\frac{\pi k}{p})$ terms of $F(p/2)$ are exactly minus those of $F(2/p)$ and cancel in the
sum; combined with Theorem~\ref{thm:F2q} this yields $F(p/2)$ in closed form.
\end{theorem}

\subsection{The derivative law}
\begin{theorem}[Radchenko--Zagier Proposition 4, $P=2$]\label{thm:deriv}
For all $q\ge3$, with $T(q):=\tfrac2q F'(2/q)-(1+\log2)$,
\[
  T(q)=\frac{q^2-4}{24q}\pi^2+W(q),\qquad
  W(q)=\sum_{k=1}^{q-1}\cot\tfrac{2\pi k}{q}\,\Cl\!\Bigl(\tfrac{2\pi k}q\Bigr),
\]
with an extra $-\log2$ for even $q$. Unlike $F(2/q)$ itself, this does \emph{not} reduce to elementary form:
$W(q)$ is a genuine Clausen combination.
\end{theorem}

\subsection{The collapse criterion and the values $J(p/q)$}
\begin{proposition}[Collapse criterion]\label{prop:collapse}
For coprime $p/q$, $F(p/q)-F(1)$ reduces to dilogarithms at \emph{rational} arguments (together with $\pi^2$
and logarithms) if and only if $\rad(q)\mid 30$, i.e.\ every prime factor of $q$ lies in $\{2,3,5\}$;
otherwise it requires genuine $\Li_2$ at the cyclotomic-ratio arguments $(1-\zeta_p^a)/(1-\zeta_q^b)$. Across
all rational cases only two new dilogarithm constants ever occur: $\Li_2(1/3)$ (conductor $3$) and the pair
$\{\Li_2(1/5),\Li_2(2/5)\}$ (conductor $5$).
\end{proposition}
\noindent Sample closed forms:
$F(3/4)-F(1)=-\tfrac{19}{144}\pi^2-\tfrac34\log^2 2-2\log2\log3+\tfrac74\log^2 3+2\Li_2(\tfrac13)$ and
$F(4/5)-F(1)=\tfrac{3}{80}\pi^2+\tfrac{11}4\log^2 2-\tfrac12\log^2 5-\Li_2(\tfrac15)-\log^2(2\sin\tfrac\pi5)-\log^2(2\sin\tfrac{2\pi}5)$.

\begin{proposition}[$J$-values]\label{prop:J}
With $J$ as in~\eqref{eq:Jint}, $J(2/q)$ is a pure-logarithm value if and only if $q\in\{3,5\}$:
$J(2/3)=-\tfrac{\pi^2}{144}+\tfrac34\log^2 2$ and
$J(2/5)=\tfrac{11\pi^2}{240}+\tfrac34\log^2 2-2\log^2\tfrac{\sqrt5+1}2$. So the paper's $J(2/5)$ --- its only
worked $J$ example --- is a special case, not a general law; the dilogarithms cancel only at $q=3,5$. The
values $J(1/q)$ always retain the $\Li_2(-\cot^2)$ dilogarithms inherited from $F(2/q)$.
\end{proposition}

\section{Bridges to the golden ladders and the Clausen lattice}\label{sec:bridges}
The new constants are not isolated; they touch two corpora we have already mapped.

\paragraph{The golden bridge ($q=5$).} The arguments $-\cot^2\tfrac\pi5=-\tfrac{5+2\sqrt5}5$ and
$-\cot^2\tfrac{2\pi}5=-\tfrac{5-2\sqrt5}5$ lie in $\mathbb Q(\sqrt5)$, the field of the golden polylogarithm
ladders. Comparing Theorem~\ref{thm:F2q} ($q=5$) with eq.~24 gives the functional equation
\[
  \Li_2\!\Bigl(\tfrac45\Bigr)+\Li_2\!\Bigl(-\cot^2\tfrac{\pi}{5}\Bigr)+\Li_2\!\Bigl(-\cot^2\tfrac{2\pi}{5}\Bigr)
  =2\log^2 2-2\log2\log\tfrac25-6\bigl(\log^2(2\sin\tfrac\pi5)+\log^2(2\sin\tfrac{2\pi}5)\bigr),
\]
hence the golden-sum closed form (with $\varphi=\tfrac{1+\sqrt5}2$)
\[
  \Li_2\!\Bigl(-\cot^2\tfrac{\pi}{5}\Bigr)+\Li_2\!\Bigl(-\cot^2\tfrac{2\pi}{5}\Bigr)
  =\Li_2\!\Bigl(\tfrac15\Bigr)-\frac{\pi^2}{6}-3\log^2\varphi+\tfrac14\log^2 5 .
\]
Only the symmetric sum reduces: the antisymmetric part
$X=\tfrac12\bigl(\Li_2(-\cot^2\tfrac\pi5)-\Li_2(-\cot^2\tfrac{2\pi}5)\bigr)$ admits no integer relation over
$\{\pi^2,\log^2\varphi,\log^2 5,\log^2 2,\Li_2(\tfrac15),\Li_2(\tfrac25),\text{log products}\}$, so it is a
genuinely new $\mathbb Q(\sqrt5)$ real-dilogarithm constant. (Consistently, $X$ is not balanced as a
Bloch-group element, so the vanishing of $K_3(\mathbb Q(\sqrt5))\otimes\mathbb Q$ does not force it to reduce.)

\paragraph{The Clausen bridge (the derivative law).} The Clausen sum $W(q)$ of Theorem~\ref{thm:deriv} is a
$\mathbb Q(\zeta_q)^+$-linear combination of the conductor-$q$ Clausen values $\Cl(\tfrac{2\pi k}q)$ --- an
element of the conductor-$q$ Clausen lattice over the real cyclotomic field. The coefficients are rational
over $\mathbb Q(\zeta_q)^+$ exactly when $\cot\tfrac{2\pi k}q\in\mathbb Q(\zeta_q)^+$, in which case $W(q)$
closes cleanly:
\[
  W(8)=2\Cl\!\Bigl(\tfrac{\pi}4\Bigr)-2\Cl\!\Bigl(\tfrac{3\pi}4\Bigr),\qquad
  W(12)=\frac{22\sqrt3}{9}\Bigl(\Cl\!\Bigl(\tfrac{\pi}6\Bigr)-\Cl\!\Bigl(\tfrac{5\pi}6\Bigr)\Bigr).
\]
For $q=5,7,9,11$ the cotangents lie in a strictly larger field, so no rational reduction exists, while
lattice membership persists. Thus differentiating $F$ at a rational moves from the real-dilogarithm world of
$F(2/q)$ into the imaginary-dilogarithm (Clausen) lattice mapped at roots of unity.

\section{Verification record}\label{sec:verif}
Every displayed identity was re-checked against the exact sum~\eqref{eq:dilogsum} (and, where applicable, the
independent integral or a numerical derivative). The Theorem-1 derivation was additionally audited
step-by-step ($10^{-149}$) and Theorem~2 confirmed to $10^{-159}$ by an independent multi-agent pass.

\begin{center}\small
\begin{tabular}{lll}
\toprule
Result & test & max residual \\
\midrule
Tables 1 \& 2 (45 rows, 2 corrected) & PARI/GP, $200$ digits & $<10^{-80}$ \\
Thm~\ref{thm:F1q} $F(1/q)$ & $q=2..40$ vs \eqref{eq:dilogsum} & $1.0\times10^{-99}$ \\
Thm~\ref{thm:F2q} $F(2/q)$ & $q=3..40$ vs \eqref{eq:dilogsum} & $1.1\times10^{-99}$ \\
Thm~\ref{thm:np1} $F(n/(n+1))$ & $n=2..16$ & $8.9\times10^{-160}$ \\
Thm~\ref{thm:p2} $F(p/2)$ reciprocity & odd $p=3..21$ & $4.1\times10^{-159}$ \\
Thm~\ref{thm:deriv} $F'(2/q)$ & $q=3..30$; vs numeric deriv. & $1.1\times10^{-129}$ \\
Prop~\ref{prop:collapse} samples $F(3/4),F(4/5)$ & PSLQ + canary & $1.6\times10^{-61}$ \\
Prop~\ref{prop:J} $J(2/3),J(2/5)$ & vs \eqref{eq:dilogsum} & $7.8\times10^{-62}$ \\
\S\ref{sec:bridges} golden bridge / sum & vs \eqref{eq:dilogsum} & $1.6\times10^{-61}$ \\
\S\ref{sec:bridges} $W(8),W(12)$ & vs definition & $1.6\times10^{-60}$ \\
\bottomrule
\end{tabular}
\end{center}

\noindent All PSLQ relations carried a Champernowne canary with coefficient $0$ (no false positives). The
reproducers are: \texttt{rational\_harvest.py} (the exact sum and the proven formulas
\texttt{F2q,F1q,Fn\_over\_np1,Fp2,dT2q}), \texttt{verify\_full.gp} (the $45$-row table check),
\texttt{audit\_thm1.py} (the step-by-step proof audit), and \texttt{bridges.py} (\S\ref{sec:bridges}), all
under \texttt{src/PolyLog/tools/scratch/herglotz/}. The verified statements are recorded in
\texttt{src/PolyLog/identities/herglotz-rational-values.wl} and
\texttt{herglotz-quadratic-units.wl}.

\section*{Summary}
Starting from a paper that records the values of $F$ at quadratic units (in two tables) and at rationals
(through two worked examples), we have: independently re-verified the tables and caught two of their typos;
promoted the two rational examples to a family of general theorems with proofs; and connected the resulting
constants to the golden ladder field $\mathbb Q(\sqrt5)$ and to the Clausen lattice at roots of unity. The
Herglotz function is now a fully verified, well-connected part of the polylogarithm corpus.

\end{document}
